% Generated draft: requires mathematical and source review.
\documentclass[11pt]{article}
\usepackage[margin=25mm]{geometry}
\usepackage{fontspec}
\setmainfont{Latin Modern Roman}
\usepackage{amsmath,amssymb,mathtools,unicode-math}
\setmathfont{Latin Modern Math}
\usepackage{enumitem,needspace,xurl,hyperref}
\setlength{\parindent}{0pt}
\setlength{\parskip}{5pt}
\setlength{\emergencystretch}{3em}
\newcommand{\N}{\mathbb{N}}
\newcommand{\Z}{\mathbb{Z}}
\newcommand{\Q}{\mathbb{Q}}
\newcommand{\R}{\mathbb{R}}
\newcommand{\C}{\mathbb{C}}
\newcommand{\F}{\mathbb{F}}
\newcommand{\A}{\mathbb{A}}
\newcommand{\PP}{\mathbb{P}}
\newcommand{\E}{\mathbb{E}}
\newcommand{\eps}{\varepsilon}
\newcommand{\dd}{\,\mathrm{d}}
\begin{document}
\section*{UnsolvedMath ID: 30002884}
\textbf{Source problem code:} catalogue entry
\subsection{Definitions and mathematical statement}
The supplied record asks whether every injective-reachable quantum cluster algebra admits a unique common triangular basis parametrized by tropical points. It associates this prediction with the Fock--Goncharov conjecture, without specifying a precise version of that conjecture.

For a quantum cluster algebra \(A\), let \(R_A\) denote its intended coefficient ring. An \(R_A\)-basis is a subset \(\mathcal B\subseteq A\) such that every element of \(A\) has a unique expression as a finite \(R_A\)-linear combination of elements of \(\mathcal B\). Let \(T_A\) denote the intended set of tropical parameters. A parametrization by \(T_A\) means a bijection \(T_A\to\mathcal B\).

The proposed assertion can therefore be displayed schematically as
\[
\forall A\in\mathscr C_{\mathrm{ir}},\quad
\exists!\,\mathcal B\subseteq A\quad
\left[
\begin{gathered}
\mathcal B\text{ is an }R_A\text{-basis of }A,\\
\mathcal B\text{ is common triangular},\\
\text{there exists a bijection }T_A\longrightarrow\mathcal B.
\end{gathered}
\right],
\]
where \(\mathscr C_{\mathrm{ir}}\) is intended to be the class of injective-reachable quantum cluster algebras under the source's conventions. Here uniqueness concerns the basis as a subset of \(A\); the record does not specify whether it also asserts a unique or canonical parametrization.

This is a schematic statement requiring further source clarification. The record does not define injective reachability or the common triangular property, specify the coefficient ring and quantum parameter, state the allowed exchange data and frozen-variable conventions, or identify the tropical parameter space and its seed transition maps. In particular, it does not supply the normalization, ordering, and seedwise conditions needed to turn ``common triangular'' into a precise predicate. These choices are left unresolved, rather than supplied as additional hypotheses. The record also does not distinguish the ordinary quantum cluster algebra from the quantum upper cluster algebra. Consequently, the supplied evidence does not determine a fully specified mathematical conjecture in the claimed generality.
\subsection{Short English statement}
The proposed conjecture says that every injective-reachable quantum cluster algebra has exactly one common triangular basis, with basis elements indexed by tropical points. The record leaves essential conventions and definitions unspecified, so the displayed formulation is schematic and requires human mathematical review. Its supplied partially solved status is dated curation, not an independently verified conclusion.
\subsection{Sources}
\begin{itemize}
\item Supplied local UnsolvedMath registry record, permanent ID \(30002884\), title: Common Triangular Bases for Quantum Cluster Algebras. Collection: Oberwolfach Reports Open Problems. Exact source problem code: \texttt{catalogue entry}. The record gives proposed year \(2015\) and names no proposer.
\item Supplied report attribution: \emph{Enveloping Algebras and Geometric Representation Theory} (2016), Oberwolfach Reports, DOI 10.4171/owr/2015/25, pp.\ 1385--1447. \url{https://doi.org/10.4171/owr/2015/25}. The supplied background describes Oberwolfach Reports as documenting mathematical workshops and their open-problem discussions. No exact problem page or passage is supplied.
\item Supplied literature source: Kang--Kashiwara--Kim--Oh, \emph{Monoidal categorification}. \url{https://doi.org/10.1090/jams/895}.
\item Supplied literature source: Gross--Hacking--Keel--Kontsevich, \emph{Canonical bases}. \url{https://arxiv.org/abs/1106.4971}.
\item Literature assessment dated exactly \texttt{2026-08-21}: the supplied status is \texttt{partially\_solved}. The supplied assessment states that common triangular-basis and monoidal-categorification statements are established for important classes, including the cited quantum-group settings, but not for the full generality of this record. This is supplied dated curation derived from the OWR material, not independently verified here; the source itself requests independent verification of time-sensitive status and cited claims.
\item The registry labels the statement \texttt{corrected\_verified} and reports a manual semantic, typographic, and LaTeX review concluding that its cleanup is meaning-preserving, self-contained, and adds no unsupported source-dependent content. This is a supplied editorial assessment, not mathematical verification of this draft or resolution of the missing conventions.
\end{itemize}
\end{document}
